%! Author = gabry
%! Date = 03/07/2026

\chapter{Introduzione}\label{ch:introduzione}

\section{Obiettivo}\label{sec:obiettivo}
Il progetto di tesi si propone di sviluppare un'applicazione mobile cross-platform, compatibile sia su piattaforme
Android che iOS, partendo dallo sviluppo di un'applicazione nativa scritta in Android Kotlin, che necessità
di una migrazione verso Kotlin Multiplatform.
L'applicazione nativa offre analisi tramite l'input di parametri dei tumori di tipo Melanoma, con metrica Melanoma TNM,
seguendo il manuale AJCC 8th Edition.
Questa sarà successivamente estesa con l'implementazione di analisi di tipo cSCC Testa-Collo, con metrica cSCC TNM, sempre seguendo
il precedente manuale AJCC 8th Edition\@.

L'obiettivo è quello di offrire uno strumento semplice e veloce per semplificare il lavoro dei medici nella fase di analisi dei tumori,
con dati sincronizzati in ambiente cloud tramite app installabili su dispositivi mobili.

Il processo di realizzazione del software è strutturata in una sequenza lineare di fasi, nello specifico:

\begin{itemize}
    \item Analisi dei requisiti;
    \item Progettazione;
    \item Sviluppo;
    \item Conclusioni;
\end{itemize}
